% BHU PYQ -- Transcribed & Corrected
% Paper: CHESE31 / CHESE21: Synthetic Methodologies (Mid-Term)
% Exam: B.Sc. (Hons) Semester III Mid-Semester Examination 2025-26 (NEP SEC)
% Department: Chemistry

\documentclass[11pt,a4paper]{article}
\usepackage[a4paper,top=2.5cm,bottom=2.5cm,left=2.8cm,right=2.8cm,headheight=14pt]{geometry}
\usepackage{amsmath,amssymb,chemformula}
\usepackage[T1]{fontenc}
\usepackage[utf8]{inputenc}
\usepackage{lmodern,microtype,fancyhdr,enumitem,hyperref,lastpage,parskip}

\hypersetup{
    colorlinks=true,
    urlcolor=black,
    linkcolor=black,
    pdftitle={CHESE31: Synthetic Methodologies -- Mid-Term 2025-26}
}

\pagestyle{fancy}
\fancyhf{}
\renewcommand{\headrulewidth}{0.4pt}
\renewcommand{\footrulewidth}{0.4pt}
\fancyhead[L]{\small   \textit{Banaras Hindu University}}
\fancyhead[R]{\small   \textit{CHESE31: Synthetic Methodologies (Mid-Term)}}
\fancyfoot[C]{\small Page  \thepage\ of \pageref{LastPage}}


\newlist{parts}{enumerate}{2}
\setlist[parts,1]{label=(\alph*),leftmargin=*,itemsep=4pt,topsep=2pt}

\newcommand{\pts}[1]{\hfill   \textit{[#1]}}


\begin{document}


\begin{center}
  {\large\bfseries BANARAS HINDU UNIVERSITY}\\[2pt]
  {\normalsize B.Sc. (Hons) Semester III Mid-Semester Examination 2025-26 (NEP SEC)}\\[6pt]
  \rule{0.8\linewidth}{0.5pt}\\[6pt]
  {\large\bfseries SEC --- CHEMISTRY}\\[4pt]
  {\large\bfseries Paper Code: CHESE31 / CHESE21 : Synthetic Methodologies}\\[4pt]
  {\normalsize Time: 1 Hour\hfill Maximum Marks: 20}
\end{center}

\rule{\linewidth}{0.4pt}

\bigskip


\begin{enumerate}[label=   \textbf{\arabic*.}]
    \item Answer the following:
    
\begin{parts}
        \item Describe the basic principle of distillation of organic solvents with suitable equations and examples. Draw the schematic diagram of solvent distillation assembly used in the chemical laboratory.\pts{4}
        \item What are molecular sieves? What are their approximate chemical compositions? Write down the classification of molecular sieves with specific examples of their use.\pts{4}
        \item Describe the drying methodologies for THF and Methanol in detail.\pts{2}
    \end{parts}

    \bigskip

    \item Answer the following:
    
\begin{parts}
        \item What are the key characteristics that define a click reaction? Explain the mechanism of the Copper(I)-catalyzed azide--alkyne cycloaddition (CuAAC) reaction.\pts{4}
        \item Describe the thiol--ene click reaction and its mechanism.\pts{3}
        \item Predict the correct products for the following reactions:\pts{3}
        
\begin{enumerate}[label=(\i)]
            \item Phenyl azide + Phenylacetylene $\xrightarrow{ \text{Cu(I)}}$
            \item 2-Methylbut-3-yn-2-ol + Phenyl azide $\xrightarrow{ \text{Cp*RuCl(PPh}_3 \text{)}_2}$
        \end{enumerate}
    \end{parts}
\end{enumerate}

\end{document}
